\documentclass[11pt]{article}

\usepackage[margin=1.05in]{geometry}
\usepackage{amsmath,amssymb,booktabs,array}
\usepackage{microtype}
\usepackage{xcolor}
\usepackage[colorlinks=true,linkcolor=blue!55!black,
            citecolor=blue!55!black,urlcolor=blue!55!black]{hyperref}
\hypersetup{
  pdftitle={Ghosts, Geometry, and Reality in Fourth-Order Gravity},
  pdfauthor={GPT-5.6.sol; Claude Fable 5; Asger Alstrup Palm},
  pdfsubject={A guided introduction to physical claims, mathematical carriers, and evidence in fourth-order gravity}
}

\newcommand{\R}{\mathbb R}
\newcommand{\C}{\mathbb C}
\newcommand{\keybox}[1]{%
  \begin{center}
  \fbox{\parbox{0.88\linewidth}{\centering #1}}
  \end{center}}

\title{Ghosts, Geometry, and Reality in Fourth-Order Gravity\\[4pt]
\large A guided introduction to physical claims, mathematical carriers,
and evidence}
\author{GPT-5.6.sol\thanks{OpenAI model and principal programme author.}
\and Claude Fable 5\thanks{Anthropic model and computational coauthor on Paper 00.}
\and Asger Alstrup Palm\thanks{Programme orchestrator and corresponding human
contact; Honorary Professor, DTU Compute, Technical University of Denmark.
\texttt{asger@area9.dk}.}}
\date{17 August 2026}

\begin{document}
\maketitle

\begin{abstract}
Fourth-order gravity contains more local solutions than general relativity.
Some appear with an indefinite sign in familiar decompositions and are called
ghosts.  The sign is a serious warning, but it does not identify by itself a
physical particle.  Local solution space, gauge-reduced classical phase
space, nonlinear continuation, and positive quantum state space are different
objects.

This paper develops a way to keep those objects apart.  A physical conclusion
is treated as a judgement depending on physical postulates, a mathematical
carrier, logical and existence assumptions, and explicit bridges to
observation.  Pais--Uhlenbeck systems, compact pure-Weyl laboratories,
Schwarzschild perturbations, Mannheim rotation curves, and the
Bateman--Turok Euclidean lattice serve as controlled examples.  Each answers
a different question; none is silently promoted into a complete theory.

The examples yield exact survival results and exact obstructions.  Additional
classical pure-Weyl directions can survive gauge reduction.  A repeated
Schwarzschild spin-two block can form a non-split extension with a defective
resonance.  Strict pure Weyl has a nonzero local Euclidean one-loop BV anomaly
in the declared complex.  A full-complex BRST--Hadamard two-point pseudo-state
can satisfy causal and gauge conditions while remaining indefinite.  A
bounded galaxy comparison can favor different models under different scores,
and an explicit positive Euclidean family can refute a proposed universal
stability shortcut.

The conclusion is methodological rather than triumphalist: the word
``ghost'' should be replaced by a ladder of falsifiable obligations, and each
obligation should state which physics, mathematics, logic, and evidence make
it true.
\end{abstract}

\section*{Authorship and how to read this paper}

Asger Alstrup Palm, a computer scientist and Honorary Professor at DTU
Compute, orchestrates the programme and is the accountable human contact.
AI systems perform substantial derivation, programming, literature work,
adversarial review, and drafting.  Claude Fable~5 contributed to this
introductory paper.  Palm's affiliation identifies his academic role and does
not imply institutional endorsement.

AI authorship is neither evidence for nor evidence against a scientific
claim.  The relevant questions are whether the assumptions are declared,
whether the derivation is valid, whether the computation can reject damaged
inputs, and whether another researcher can locate the exact boundary of the
claim.  Sources, certificates, verifiers, and papers are available in the
\href{https://github.com/area9innovation/weyl-gravity}{public repository}.

\tableofcontents

\section{An equation is not yet a theory}

An elegant field equation does not say by itself what exists, what can be
prepared, what evolves causally, or what an observer will count.  A usable
theory requires at least five layers:
\begin{enumerate}
\item \textbf{Physical postulates:} symmetry, locality, causality,
conservation, and composition rules.
\item \textbf{A mathematical carrier:} finite algebra, Hilbert or Krein
space, operator algebra, smooth fields, distributions, or an internal
geometric setting.
\item \textbf{Logic and existence assumptions:} classical or constructive
logic, forms of infinity, completeness, compactness, and choice.
\item \textbf{Operational bridges:} states, observables, probability,
preparation, clocks, detectors, and reconstruction.
\item \textbf{Empirical contact:} a numerical prediction with declared inputs,
uncertainties, and a falsifiable comparison protocol.
\end{enumerate}

The programme writes this dependence schematically as
\[
 L+S+M+\operatorname{Enc}(P)\ \vdash\ O,
\]
where \(L\) is the logic, \(S\) the set or existence theory, \(M\) the
mathematical carrier, \(P\) the physical postulates, and \(O\) the obligation
being proved.  The notation is not a claim that physics reduces to formal
logic.  It is a bookkeeping rule: if the carrier or existence principles
change, the conclusion must be re-established rather than inherited by
analogy.

\keybox{\textbf{Central rule.}
A finite identity is not a continuum limit.  A Euclidean measure is not a
Lorentzian causal state.  A cohomology class is not a particle.  A fitted
curve is not a complete physical explanation.}

\section{The ghost question is four questions}

A nondegenerate fourth-order oscillator has an equation of the form
\[
 \left(\frac{d^2}{dt^2}+\omega_1^2\right)
 \left(\frac{d^2}{dt^2}+\omega_2^2\right)z=0.
\]
It contains two second-order oscillatory branches.  In the standard
Ostrogradsky description, one normally appears with the opposite Hamiltonian
sign~\cite{Ostrogradsky,Woodard,PU1950}.  Bender and Mannheim showed that the
unequal-frequency Pais--Uhlenbeck model also admits a non-Hermitian
PT-symmetric description with a positive metric~\cite{BM2008PRL,BM2008PRD}.
At equal frequency the branches coalesce into a Jordan system.

None of these facts settles gravity.  In a gauge field theory the word
``ghost'' is commonly asked to carry four distinct meanings:
\begin{enumerate}
\item a local solution of a differential equation;
\item a classical direction surviving gauge symmetry, constraints, charges,
boundaries, and the symplectic radical;
\item a direction admitting consistent nonlinear continuation;
\item a state or particle in a positive quantum theory.
\end{enumerate}

The implications between these levels are neither automatic nor symmetric.
A negative propagator residue warns of a problem but is not yet a
gauge-reduced state.  A nonzero classical symplectic pairing proves that a
direction is not radical, but does not provide a Born rule.  A positive metric
for one free block does not guarantee an interaction-compatible metric.  A
well-behaved causal two-point distribution need not be positive on physical
cohomology.

\section{Why use several model universes?}

No single background makes every obligation tractable.  The programme uses
declared laboratories, each chosen to expose a particular interface.

\begin{center}
\begin{tabular}{>{\raggedright\arraybackslash}p{0.30\linewidth}
                >{\raggedright\arraybackslash}p{0.60\linewidth}}
\toprule
Laboratory & Question isolated\\
\midrule
Pais--Uhlenbeck oscillator and scalar field
& Positive and Krein completions, vacuum representations, Jordan limits,
and interaction obstructions.\\
Conformal cylinder and compact product backgrounds
& Full free gauge complexes, residual cohomology, causal transport,
action-derived pairings, and nonlinear balance.\\
Schwarzschild exterior and static spherical metrics
& Horizon and radiation conditions, extension structure, resonances,
Green poles, reconstruction, and static charge normalization.\\
Mannheim--Kazanas geometry and NGC~3198
& The separation between a fourth-order radial solution, matter and frame
assumptions, fitted parameters, and an empirical score.\\
Bateman--Turok periodic Euclidean lattice
& Whether a proposed all-field stability estimate survives an explicit
positive, nonseparable counterfamily.\\
Finite, constructive, Hilbert, Krein, algebraic, and internal settings
& Which physical obligations survive a change of mathematical foundation or
carrier.\\
\bottomrule
\end{tabular}
\end{center}

The laboratories are not asserted to be one universe.  They are controlled
experiments on the reasoning.  A compact positive block does not establish an
asymptotically flat particle; a Euclidean lattice does not establish
Lorentzian reconstruction; and a black-hole endpoint theorem does not
establish cosmology.

\section{From oscillator signs to interacting fields}

The elementary fourth-order models distinguish several structures often
compressed into ``the metric'': a finite-dimensional similarity transform, a
Krein real form, a vacuum covariance, a field representation, a Fock sector,
and the equal-frequency Jordan limit.  These structures obey different
continuity and ultraviolet requirements.

Linearized Einstein--Weyl gravity adds Lorentz covariance, helicity, gauge
constraints, and a conformal critical limit.  Gauge reduction removes
redundant variables but does not turn the gravitational problem into
independent copies of a scalar oscillator.  The real form must be compatible
with the tensorial and gauge structure.

Interaction is a separate gate.  In the scalar and gravitational channels
tested, resonant branch conversion obstructs the canonical analytic
continuation of the free positive metric.  The conclusion is scoped: it does
not exclude every indefinite, nonlocal, singular, or nonperturbative
prescription.  It proves that a free completion cannot be advertised as an
interacting solution without another theorem.

This distinction is especially important in higher-derivative gravity.
Stelle's renormalizability result~\cite{Stelle} concerns ultraviolet power
counting and counterterms; it does not by itself settle the physical state
space.  Conversely, a difficulty with one state-space completion does not
erase the classical field equations or their gauge cohomology.

\section{Pure Weyl gravity across the obligation ladder}

\subsection{Residual cohomology is not a graviton count}

On the conformal cylinder, a declared zero-charge reduction treats all
fifteen residual conformal generators as constraints.  The first surviving
centered classes are
\[
 H^4_{\rm res}
 =\operatorname{span}\{[W_+^2],[W_-^2]\},
 \qquad G_{\rm res}=I_2.
\]
They are deformation or vertex classes, not one-particle gravitons.  A
companion causal BV--BFV construction transports the relevant free pairings
between support conditions~\cite{BVBFV}.  Neither statement supplies a
positive interacting particle space.

Compact axial and polar laboratories show that some additional fourth-order
directions survive gauge reduction and have nonzero action-derived Lee--Wald
pairings.  Their sign pattern is not the shortcut ``Einstein positive,
additional negative.''  Other directions fail global charge balance,
secular-growth, resonance, or Hamiltonian health tests.  The scientifically
stable object is therefore a classification of promotion and obstruction,
not a universal rescue or universal no-go theorem.

\subsection{The classical complex is an input to quantization}

The classical BV--BFV complex is imported into the quantum programme through
an immutable, content-addressed snapshot.  Independent checks cover the
declared nilpotency, contraction, chain-map, cyclicity, residual-cohomology,
and pairing obligations.  This gate matters because rebuilding an unofficial
classical complex inside the quantum calculation could make a later quantum
identity true of the wrong theory.

Typed \(q_2/q_3\) compatibility with advanced and retarded Green homotopies
then connects the nonlinear brackets to the causal classical data.  The
construction does not supply arbitrary all-order response trees, convergence,
or an effective global Green solver.

\subsection{Anomaly and pseudo-state are different quantum questions}

In the declared local Euclidean quotient, strict fixed-field-content pure
Weyl gravity has a nonzero one-loop BV anomaly.  A formal shifting compensator
makes the tested cocycle exact only after enlarging the field and counterterm
space.  That is a restoration in a changed theory, not proof that strict pure
Weyl is anomaly-free.

On the Lorentzian side, the full 386-row off-shell carrier admits a
BRST--Hadamard two-point pseudo-state pair satisfying the declared bisolution,
graded commutator, wavefront, Ward, reality, and stationarity conditions.  The
word \emph{pseudo-state} is essential.  The pair is indefinite and does not
establish positivity on physical cohomology, a positive graviton Hilbert
space, a Feynman propagator, renormalized time-ordered products, a Lorentzian
quantum master equation, scattering, or unitarity.

The anomaly and pseudo-state results do not repair one another.  They answer
different obligations on different dependency rails.

\section{Black holes: boundary selection and defective response}

Maldacena-type conditions can isolate an Einstein sector of conformal gravity
in Euclidean anti-de Sitter or late-time de Sitter settings~\cite{Maldacena}.
The Lorentzian Schwarzschild question is different: do future-horizon
regularity and ordinary one-sided radiative traces force the additional
pure-Weyl curvature carrier to vanish?

For the complete axial
\(\ell=2\) Schwarzschild system, the answer is no.  Its radial differential
module has successive factors
\[
 M_{\rm RW}^{(2)},\qquad M_{\rm RW}^{(2)},\qquad M_{\rm RW}^{(1)}.
\]
The repeated spin-two factors form a non-split self-extension rather than a
direct sum.  Exact Wronskian arguments populate every incoming direction by a
future-horizon-regular exterior solution for real positive frequency.  The
same filtration excludes exponentially growing separated axial modes.

This coexistence is instructive.  Separated-mode stability says that one
particular instability is absent.  It does not make the inherited flux
positive and does not bound the full time-domain resolvent.

At one validated simple Regge--Wheeler quasinormal frequency~\cite{ReggeWheeler,Leaver},
the complete connection has Smith valuations
\[
 (0,0,2).
\]
The algebraic multiplicity is two, the geometric multiplicity is one, and the
root chain has length two.  The geometric root lies in the Einstein
submodule; the generalized root has nonzero gauge-invariant traceless-Ricci
carrier.  This is related to the logarithmic-partner structure of critical
gravity~\cite{LuPope,LogGravity}, although the asymptotically flat endpoint
tangent is polynomial rather than an independent literal \(\log r\) mode.

The compactly cut-off exterior radial Green operator has
\[
 R(\omega)
 =\frac{R_{-2}}{(\omega-\omega_n)^2}
  +\frac{R_{-1}}{\omega-\omega_n}
  +R_{\rm hol}(\omega),
 \qquad R_{-2}\ne0.
\]
The same pole appears in a fixed-domain two-ended complex-scaled Fredholm
realization and is the meromorphic continuation of a retarded, mode-reduced,
compact-source/compact-observation transfer.  An isolated local contour
therefore contains \(e^{i\omega_n t}(V_1+i tV_0)\).

\keybox{\textbf{Boundary of the resonance claim.}
The non-split extension, double radial Green pole, and retarded mode-reduced
parent are established in their declared settings.  A full metric/BV
retarded propagator, global inverse-Laplace deformation, complete late-time
expansion, and real astrophysical detector waveform are not.}

Static Bach-flat black holes answer another question.  Residual-basic
normalization yields a consistent Hamiltonian, Wald entropy, and signed
temperature on the certified static family, but this is not a radiative or
physical-process thermodynamics theorem.

\section{Two tests beyond black-hole perturbation theory}

\subsection{Mannheim gravity: solution, interpretation, and fit}

The Mannheim--Kazanas geometry contains a linear radial potential because the
static fourth-order vacuum equation admits it~\cite{MK1989,MannheimOBrien2012}.
That mathematical fact does not determine how the term is sourced or how
massive matter moves.  A physical explanation must also specify the matter
sector, scale breaking, conformal frame, source moments, and global boundary
data.

The distinction becomes visible in a bounded 39-point NGC~3198 protocol:
\begin{center}
\begin{tabular}{lrrc}
\toprule
Model & RMS (km/s) & Reduced \(\chi^2\) & Random-error gate\\
\midrule
Newtonian baryons & 23.896 & 128.722 & fail\\
GR plus NFW halo & 5.148 & 0.965 & pass\\
Mannheim conformal gravity & 4.694 & 3.202 & fail\\
\bottomrule
\end{tabular}
\end{center}

Mannheim has the smallest unweighted residual, whereas GR+NFW has the better
uncertainty-weighted score.  This is not a contradiction: the two scores ask
different questions.  Nor is it a population-level verdict.  Distance,
inclination, beam smearing, stellar populations, gas profiles, and other
systematics lie outside the frozen protocol.

\subsection{Bateman--Turok: a universal shortcut meets a counterfamily}

Bateman and Turok propose a positive Euclidean lattice and a hidden ghost
parity as a route around the usual higher-derivative instability
story~\cite{BatemanTurok2026}.  One tempting deterministic strategy is to
bound the complete action gradient below at the free bilaplacian scale for
every positive field.

An explicit periodic-torus family disproves that universal bound.  On
\(L_n=n^{24}\), an exactly mean-compensated critical source is inverted by the
periodic graph Green operator, lifted above its minimum, and coupled weakly in
four directions.  The fields remain positive, nonseparable, and of
polynomial contrast.  Their action stays bounded above and away from zero,
while
\[
 \frac{Q_n}{\omega_{L_n}^2}=O(n^{-2})\longrightarrow0.
\]

This refutes one deterministic all-field scaled-PL architecture, not the full
Bateman--Turok programme.  Gibbs typicality, the interacting Witten quotient,
the \(H^{-1}\) moment, the continuum measure, reflection positivity, and
Lorentzian reconstruction remain logically separate tests.

\section{Reverse foundations: changing the mathematical world}

The previous examples change physical models while mostly retaining familiar
mathematics.  Reverse foundations asks the complementary question: what
happens when the physical obligation is held fixed but the mathematical or
logical background changes?  This extends the idea of reverse physics---work
backward from what a law must accomplish to the assumptions that make it
possible~\cite{CarcassiAidala2022}.

The interactive atlas crosses three axes:
\begin{itemize}
\item six foundational regimes, including classical, weakened-choice,
constructive, internal/topos, and finite settings;
\item six carrier families, including finite algebra, Hilbert space, Krein
space, operator algebra, smooth/distributional fields, and localic geometry;
\item sixteen physical obligations, from state existence and dynamics through
causality, gauge reduction, interaction, renormalization, reconstruction, and
empirical adequacy.
\end{itemize}

The Cartesian product contains \(6\times6\times16=576\) coordinates.  They
are not 576 candidate universes.  They are 576 scoped questions.  A populated
cell records evidence about one obligation in one setting; it does not show
that neighboring cells compose.

Theory assemblies group evidence by research tradition, while theory
passports follow a single route through foundations, state space, dynamics,
observable, prediction, and empirical benchmark.  This exposes an important
difference between \emph{coverage} and \emph{completion}.  A mature tradition
may possess tools for every job without those tools belonging to one model.
An unconventional programme may solve a difficult late-stage identity while
lacking an earlier state or observable bridge.

Finite and constructive examples can avoid particular basis-selection or
completion assumptions.  Lean/Physlib proof passports can replay finite
identities with explicit types and signs.  Such results show that an
assumption is avoidable in that representation.  They do not prove that the
universe is finite, that all physics is constructive, or that formal proof
supplies a missing physical premise.

\section{Evidence is part of the theory interface}

To prevent scope changes from becoming invisible, quantum results carry one
or more exact dependency tags:
\begin{center}
\texttt{LOCAL-ALGEBRAIC}\quad
\texttt{EUCLIDEAN-SPECTRAL}\quad
\texttt{REDUCED-MODE}\quad
\texttt{LORENTZIAN-CAUSAL}.
\end{center}
A result does not move between these rails without an explicit bridge.
Likewise, anomaly classification, coefficient computation, QME restoration,
residual transfer, and Lorentzian certification are distinct lifecycle
states.

The preferred audit chain is
\[
 \text{claim}
 \longrightarrow \text{producer}
 \longrightarrow \text{certificate}
 \longrightarrow \text{independent verifier}
 \longrightarrow \text{adversarial test}
 \longrightarrow \text{paper}.
\]
Re-running the producer is reproduction, not independent verification.  Exact
algebra and numerical validation remain different evidence types.  A timeout,
missing input, or skipped exhaustive calculation is not a pass.

This architecture is especially important for AI-orchestrated research.  AI
systems can derive equations, write software, search literature, and produce
fluent exposition.  They can also share hidden assumptions, invent
references, and produce mutually reinforcing errors.  The answer is not to
treat AI as an oracle, but to make every substantial claim easier to attack.
A reproducible refutation is scientifically useful because it removes an
invalid route and sharpens the surviving alternatives.

\section{What the examples establish}

Six conclusions cut across the different laboratories:
\begin{enumerate}
\item \textbf{A sign is not a state.}  Pole residues, classical fluxes,
Krein signatures, and quantum probabilities are different objects.
\item \textbf{Reduction can preserve or remove extra directions.}  Gauge,
charges, boundaries, and the symplectic radical must be tested explicitly.
\item \textbf{Interactions are promotion gates.}  A free positive completion
can fail when branches mix, while a formal nonlinear solution can fail global
balance or boundedness.
\item \textbf{Causality and positivity are independent obligations.}  A
BRST--Hadamard pseudo-state can have the required singularity and gauge
identities without defining a positive physical state.
\item \textbf{A counterexample should be scoped as carefully as a theorem.}
The BT family refutes one all-field bound; the NGC~3198 gate tests one frozen
data protocol; neither decides an entire research tradition.
\item \textbf{Changing mathematics changes proof obligations.}  Replacing a
Hilbert carrier, weakening Choice, or restricting to a finite model can remove
an assumption, but operational meaning must be rebuilt rather than presumed.
\end{enumerate}

These conclusions explain why collecting individually attractive ingredients
is not enough.  A physical theory requires typed interfaces between its
foundations, states, dynamics, observables, predictions, and empirical tests.
The interfaces are often where the decisive theorem or obstruction lives.

\section{Routes into the technical evidence}

The \href{../foundations/site/index.html}{Reverse Physics Atlas} provides the
interactive matrix, dimension guide, programme assemblies, theory journeys,
Weyl BV routes, literature records, certificates, and formal-proof passports.
The companion papers can be entered thematically:
\begin{itemize}
\item \href{15-four-level-ghost-classification-phase1-synthesis.pdf}{the
four-level ghost classification} develops the separation between local,
classical, nonlinear, and quantum claims;
\item \href{19-what-conformal-gravity-must-assume.pdf}{the Mannheim analysis}
separates the radial solution from matter, frame, boundary, and empirical
premises;
\item \href{20-is-the-gauge-algebra-an-assumption.pdf}{the gauge-algebra
analysis} asks whether algebraic structure is input, derived, or
representation-dependent;
\item \href{21-reverse-foundations-of-physics.pdf}{the reverse-foundations
synthesis} develops the atlas and literature framework;
\item \href{22-bateman-turok-euclidean-torus-collapse.pdf}{the
Bateman--Turok analysis} proves the compensated Green-tail counterfamily.
\end{itemize}

These routes are thematic references, not a sequence of release notes.  Each
technical claim should be evaluated through its own assumptions, certificate,
verifier, and stated non-claims.

\section{Conclusion}

Fourth-order gravity cannot be assessed by counting poles or attaching the
word ``ghost'' to every additional solution.  The physically relevant
question is which object survives each promotion gate: local equation,
classical reduction, nonlinear continuation, causal construction, positive
quantum state, operational prediction, and empirical test.

The examples do not establish pure Weyl gravity, Mannheim conformal gravity,
or Bateman--Turok physics as a successful description of nature.  Nor do they
support a single generic dismissal.  They establish a mixed landscape of
survival theorems, obstructions, counterexamples, and unproved interfaces.

\keybox{\textbf{Durable conclusion.}
Physical postulates, mathematical carriers, logical strength, causal and
quantum claim types, interface proofs, and empirical gates must be separated.
Only explicit bridges may join them.}

That principle turns a vague dispute about ghosts into a collection of
answerable questions.  It also gives researchers a practical standard for
reviewing theories built with unfamiliar mathematics or substantial AI
assistance: inspect the obligation, inspect its dependencies, and inspect the
bridge to the next claim.

\begin{thebibliography}{99}
\raggedright
\small

\bibitem{Ostrogradsky}
M.~Ostrogradsky,
\emph{M\'emoire sur les \'equations diff\'erentielles relatives au
probl\`eme des isop\'erim\`etres},
M\'em.\ Acad.\ St.\ P\'etersbourg \textbf{VI 4} (1850), 385.

\bibitem{Woodard}
R.~P.~Woodard,
\emph{The theorem of Ostrogradsky},
Scholarpedia \textbf{10} (2015), 32243; arXiv:1506.02210.

\bibitem{PU1950}
A.~Pais and G.~E.~Uhlenbeck,
\emph{On field theories with non-localized action},
Phys.\ Rev.\ \textbf{79} (1950), 145--165,
\href{https://doi.org/10.1103/physrev.79.145}{doi:10.1103/physrev.79.145}.

\bibitem{BM2008PRL}
C.~M.~Bender and P.~D.~Mannheim,
\emph{No-ghost theorem for the fourth-order derivative
Pais--Uhlenbeck oscillator model},
Phys.\ Rev.\ Lett.\ \textbf{100} (2008), 110402,
\href{https://doi.org/10.1103/physrevlett.100.110402}{doi:10.1103/physrevlett.100.110402}.

\bibitem{BM2008PRD}
C.~M.~Bender and P.~D.~Mannheim,
\emph{Exactly solvable PT-symmetric Hamiltonian having no Hermitian
counterpart},
Phys.\ Rev.\ D \textbf{78} (2008), 025022; arXiv:0804.4190.

\bibitem{Stelle}
K.~S.~Stelle,
\emph{Renormalization of higher-derivative quantum gravity},
Phys.\ Rev.\ D \textbf{16} (1977), 953--969,
\href{https://doi.org/10.1103/physrevd.16.953}{doi:10.1103/physrevd.16.953}.

\bibitem{BVBFV}
A.~S.~Cattaneo, P.~Mnev and N.~Reshetikhin,
\emph{Classical BV theories on manifolds with boundary},
Commun.\ Math.\ Phys.\ \textbf{332} (2014), 535--603;
arXiv:1201.0290.

\bibitem{Maldacena}
J.~Maldacena,
\emph{Einstein gravity from conformal gravity},
arXiv:1105.5632.

\bibitem{ReggeWheeler}
T.~Regge and J.~A.~Wheeler,
\emph{Stability of a Schwarzschild singularity},
Phys.\ Rev.\ \textbf{108} (1957), 1063--1069,
\href{https://doi.org/10.1103/physrev.108.1063}{doi:10.1103/physrev.108.1063}.

\bibitem{Leaver}
E.~W.~Leaver,
\emph{An analytic representation for the quasinormal modes of Kerr black
holes},
Proc.\ R.\ Soc.\ Lond.\ A \textbf{402} (1985), 285--298,
\href{https://doi.org/10.1098/rspa.1985.0119}{doi:10.1098/rspa.1985.0119}.

\bibitem{LuPope}
H.~L\"u and C.~N.~Pope,
\emph{Critical gravity in four dimensions},
Phys.\ Rev.\ Lett.\ \textbf{106} (2011), 181302; arXiv:1101.1971.

\bibitem{LogGravity}
E.~A.~Bergshoeff, O.~Hohm, J.~Rosseel and P.~K.~Townsend,
\emph{Modes of log gravity},
Phys.\ Rev.\ D \textbf{83} (2011), 104038; arXiv:1102.4091.

\bibitem{MK1989}
P.~D.~Mannheim and D.~Kazanas,
\emph{Exact vacuum solution to conformal Weyl gravity and galactic rotation
curves}, Astrophys.\ J.\ \textbf{342} (1989), 635.

\bibitem{MannheimOBrien2012}
P.~D.~Mannheim and J.~G.~O'Brien,
\emph{Galactic rotation curves in conformal gravity},
J.\ Phys.\ Conf.\ Ser.\ \textbf{437} (2013), 012002; arXiv:1211.0188.

\bibitem{CarcassiAidala2022}
G.~Carcassi and C.~A.~Aidala,
\emph{Reverse Physics: From Laws to Physical Assumptions},
Found.\ Phys.\ \textbf{52} (2022), 40; arXiv:2111.09107.

\bibitem{BatemanTurok2026}
S.~Bateman and N.~Turok,
\emph{Escape from Ostrogradsky via hidden ghost parity}, 2026;
\href{https://arxiv.org/abs/2607.00096}{arXiv:2607.00096}.

\end{thebibliography}

\end{document}
